\section{Discussion}
\label{sec:discussion}

\subsection{Detection and identification answer different questions}

The detector asks whether a domain resembles anything in its training distribution. The
matcher asks which physical reference is compatible with it. Keeping them separate is what
allows a domain to be flagged as novel and then remain unidentified, which is the correct
outcome when a structure is genuinely outside the reference library. A single system that
assigned every domain to a nearest known class could not express this.

The two also fail differently, which is useful. The detector is sensitive to the descriptor
convention (\cref{sec:res-controls}) but indifferent to chemistry; the matcher is tied to the
contents of the reference library but interpretable in physical terms. A candidate accepted
by neither is a stronger signal than one rejected by either alone.

\subsection{Spatial domains are the right unit}

Correlation functions and domain-level fingerprints are only meaningful on a region that
contains one structure. Local-environment discovery alone does not define that region: it
labels particles, and the mapping from labelled particles to the domain over which
$g(r)$, $S(k)$ or a pooled descriptor should be evaluated is a separate step. Averaging
across a two-phase region produces a fingerprint belonging to neither phase.

\subsection{Abstention as a scientific outcome}

The reference-removal experiment (\cref{sec:res-abstain}) is the clearest evidence that
abstention is principled here. With the correct reference deleted, the nearest surviving
reference sits at a distance typical of two unrelated prototypes, and the system declines to
match. Reporting \emph{library-unmatched} is more informative than a confident wrong label,
because it identifies exactly the domains that warrant inspection.

\subsection{Threshold calibration is the weak point}

The score and the threshold have very different stability. Resampling the calibration groups
moves the threshold enough to change decisions, whereas the underlying distances are stable.
The predeclared rule additionally compares a group-mean statistic against individual-domain
scores, which inflates the false-positive rate to roughly a quarter of in-distribution
validation domains. Alternative rules reduce this by a factor of four or more on the same
data. This is a calibration problem rather than a representation problem, and we have
deliberately not resolved it by selecting the best-performing rule after the fact.

\subsection{Limitations}

Several limitations are explicit in the present benchmark.

\emph{Few families.} Three structural families support the detector and seven the matcher.
Perfect leave-one-family-out ranking on evaluation sets of this size does not establish
general behaviour.

\emph{Thermal ensembles are small.} Acceptance tolerances derive from eleven frames of a
single trajectory per family. These frames are correlated, so the effective sample is
smaller still.

\emph{Mean pooling dilutes minority novelty.} A domain-level mean averages a small novel
inclusion towards the familiar mean, so a nucleus of a new phase inside a familiar domain can
pass undetected.

\emph{Reference construction.} Non-cubic references were built with a single uniform lattice
scale, which preserves axial ratios appropriate to an atomic crystal rather than to a
colloidal one (\cref{sec:refbuild}). Nineteen of forty-two catalogued prototypes are
non-cubic, and one of the seven matcher references changes materially when this is corrected.

\emph{Centre-convention sensitivity.} The detector accepts domains under the frozen
reliable-centre convention and rejects the same domains under alternative centre sets. Any
comparison against externally constructed domains must use the same convention.

\emph{Near-degenerate structures.} \TODO{VERIFY: whether the matcher separates the A$_3$B
prototypes reliably, given that several share space group and stoichiometry.}

\emph{Expansion is untested.} The workflow is designed for incremental extension, but the
sequential experiment that would demonstrate it has not been run.
